% 学院 Word 模板中的参考文献著录示例及说明，保留原有颜色和顺序。
% 正式论文使用 references.bib 自动生成实际文献。
\thesisbackchapter{参考文献}
\thesisreferenceguidance
\begingroup\setlength{\parindent}{0bp}\frenchspacing\heiti\fontsize{10.56bp}{19.5bp}\selectfont
\xeCJKsetup{CJKglue=\hskip-0.06bp}
{\color{red}\noindent
{\xeCJKsetup{PunctStyle=banjiao}参考文献是毕业设计中引用文献出处的目录表。参考文献一律放在文后，书写格式要按国家标准GB7714\linebreak
－\hspace{0bp}87规定。\par}
a.\hspace{2.64bp}专著、论文集、学位论文、报告\par
[序号]\hspace{5.28bp}著者.\hspace{5.28bp}文献题名\nobreak\hspace{0bp}[文献类型标识].\hspace{5.28bp}出版地：出版者，出版年.\hspace{5.28bp}起止页码\nobreak\hspace{0bp}(任选).\par}
{\color{blue}\noindent
{\color{red}[1]}\hspace{5.28bp}刘国钧，陈绍业，王凤翥.\hspace{5.28bp}图书馆目录\nobreak\hspace{0bp}[M].\hspace{5.28bp}北京：高等教育出版社，1957. 15-18.\par}
{\color{red}\noindent
[2]\hspace{5.28bp}辛希孟.\hspace{5.28bp}信息技术与信息服务国际研讨会论文集：A集\nobreak\hspace{0bp}[C].\hspace{5.28bp}北京：中国社会科学出版社，1994.\par
[3]\hspace{5.28bp}张筑生.\hspace{5.28bp}微分半动力系统的不变集\nobreak\hspace{0bp}[D].\hspace{5.28bp}北京：北京大学数学系数学研究所，1983.\par
[4]\hspace{5.28bp}冯西桥.\hspace{5.28bp}核反应堆压力管道与压力容器的LBB分析\nobreak\hspace{0bp}[R].\hspace{5.28bp}北京：清华大学核能技术设计研究院，1997.\par
b.\hspace{5.28bp}期刊文章\par
[序号]\hspace{5.28bp}著者.\hspace{5.28bp}文献题名\nobreak\hspace{0bp}[J].\hspace{5.28bp}刊名，年，卷\nobreak\hspace{0bp}(期)：起止页码.\par
[5]\hspace{5.28bp}何龄修.\hspace{5.28bp}读顾城《南明史》[J].\hspace{5.28bp}中国史研究，1998，(3)：167-173.\par}
{\color{blue}\noindent
\xeCJKsetup{CJKglue=\hskip-0.06bp plus0.08\baselineskip}{\color{red}[6]}\hspace{5.28bp}金显贺，王昌长，王忠东，等.\hspace{5.28bp}一种用于在线检测局部放电的数字滤波技术\nobreak\hspace{0bp}[J].\hspace{0bp}清华大学学报\nobreak\hspace{0bp}(自然\linebreak
科学版)，1993，33(4)：62-67.\par}
{\color{red}\noindent
c.\hspace{5.28bp}电子文献\par
[序号]\hspace{5.28bp}著者.\hspace{5.28bp}电子文献题名\nobreak\hspace{0bp}[文献类型标识/载体类型标识].\hspace{5.28bp}电子文献的出处或可获得地址，发表或更新日期/引用日期\nobreak\hspace{0bp}(任选).\par}
{\color{blue}\noindent
\xeCJKsetup{CJKglue=\hskip2.94bp}{\color{red}[7]}\hspace{2.928bp}郭路\hspace{3bp}.\hspace{8.152bp}XML\hspace{8.328bp}数据传输的安全加密\nobreak\hspace{3.107bp}[EB/OL].\hspace{7.55bp}{\addfontfeatures{LetterSpace=-0.235}http://www-900.ibm.com/developerWorks/cn/xml/}\\
xmlb2b/index5.shtml, 2001-06/2001-10-04.\par}
\vspace{19.5bp}
{\color{red}\noindent
注1：参考文献类型\par
参考文献类型\quad 专著\quad 论文集\quad 报纸文章\quad 期刊文章\quad 学位论文\quad 报告\quad 标准\quad 专利\par
\makebox[73.6bp][l]{文献类型标识}\makebox[51.3bp][l]{M}\makebox[48.9bp][l]{C}\makebox[49.5bp][l]{N}\makebox[56.5bp][l]{J}\makebox[44.3bp][l]{D}\makebox[28bp][l]{R}\makebox[32bp][l]{S}P\par
对于其他未说明的文献类型，建议采用单字母“Z”。\par
注2：电子文献类型\par
\makebox[105.2bp][l]{电子参考文献类型}\makebox[52.5bp][l]{数据库}\makebox[73.5bp][l]{计算机程序}电子公告\par
\makebox[105.1bp][l]{电子文献类型标识}\makebox[77.5bp][l]{DB}\makebox[70.2bp][l]{CP}EB\par
注3：电子载体类型\par
\makebox[89.3bp][l]{电子文献载体类型}\makebox[26.3bp][l]{磁带}\makebox[26.3bp][l]{磁盘}\makebox[26.3bp][l]{光盘}联机网络\par
\makebox[89.4bp][l]{电子载体类型标识}\makebox[18.2bp][l]{MT}\makebox[17.7bp][l]{DK}\makebox[17.3bp][l]{CD}OL\par
如：[DB/OL]—联机网上数据库、[DB/MT]—磁带数据库、[M/CD]—光盘图书\par
[CP/DK]—\hspace{5.16bp}磁盘软件、[J/OL]—\hspace{5.16bp}网上期刊、[EB/OL]—\hspace{5.16bp}网上电子公告\par
在正文中引用参考文献：在引用处最后一个字的右上角，用方括号标明此序号。\par
如，“XML秘钥信息服务规范\nobreak\hspace{0bp}{\color{red}[1]}”、“CAST对称秘钥算法\nobreak\hspace{0bp}[2],[4]”。\par}
\endgroup
